\chapter{Sicurezza nei OS e Virtualizzazione}

\section{Introduzione alla Sicurezza dei Sistemi Operativi}
\begin{warnbox}[Vulnerabilità Intrinseche e Deployment]
È imperativo riconoscere che i Sistemi Operativi (OS) presentano vulnerabilità intrinseche, costantemente sondate da agenti malevoli (es. \textit{worm} o \textit{botnet} su Internet). Un sistema può essere compromesso già durante la delicata fase di installazione (\textit{deployment}), prima ancora che vengano applicate le patch di sicurezza o implementate le policy aziendali.
\end{warnbox}

\section{Il Processo di Messa in Sicurezza (NIST SP 800-123)}
\begin{exbox}[Le Fasi del Deployment Sicuro]
\begin{enumerate}
    \item \textbf{Valutazione dei Rischi e Pianificazione:} Analisi preventiva del contesto operativo, delle minacce e dei requisiti di business.
    \item \textbf{Hardening dell'OS base:} Messa in sicurezza del sistema operativo prima dell'esposizione in rete.
    \item \textbf{Sicurezza delle Applicazioni:} Configurazione sicura dei servizi applicativi chiave.
    \item \textbf{Protezione del Contenuto:} Salvaguardia dei dati critici tramite crittografia (Data at Rest / Data in Transit) e rigorosi controlli di accesso.
    \item \textbf{Sicurezza di Rete:} Implementazione di firewall host-based e sistemi di rilevamento delle intrusioni.
    \item \textbf{Manutenzione Continua:} Definizione di processi per il patching automatizzato, il logging e il backup nel tempo.
\end{enumerate}
\end{exbox}

\section{Hardening del Sistema Operativo}
La messa in sicurezza dell'OS rappresenta il fondamento fiduciario (\textit{Root of Trust}) per tutte le applicazioni sovrastanti. 

\begin{notebox}[Insecure by Default]
Le configurazioni di \textit{default} fornite dai produttori (sia per sistemi operativi che per dispositivi di rete) privilegiano storicamente l'usabilità immediata e la retrocompatibilità a scapito della sicurezza. Questo approccio \textit{insecure by default} richiede un'azione proattiva da parte dell'amministratore.
\end{notebox}

\begin{defbox}[Hardening]
L'hardening (indurimento) è il processo di riduzione della superficie di attacco di un sistema tramite la disattivazione di servizi non necessari, l'applicazione di configurazioni restrittive e l'implementazione del principio del minimo privilegio.
\end{defbox}

L'hardening procede attraverso i seguenti principi:
\begin{itemize}
    \item \textbf{Installazione Sicura:} Il deployment iniziale deve avvenire in una rete logicamente isolata (VLAN di \textit{staging} o rete \textit{air-gapped}) per prevenire infezioni pre-patching da parte di \textit{worm} diffusi nella rete aziendale.
    \item \textbf{Principio di Minimalità:} Installare esclusivamente i pacchetti, i moduli e i driver strettamente necessari. La rimozione (o la mancata installazione) di un servizio riduce drasticamente i vettori di attacco esposti.
    \item \textbf{Sicurezza del Boot:} Proteggere il processo di avvio (es. impostando password per UEFI/BIOS e abilitando il \textit{Secure Boot}) per prevenire l'alterazione dei file critici, il caricamento di \textit{rootkit} o di \textit{hypervisor} malevoli prima che l'OS principale venga caricato.
\end{itemize}

\subsection{Manutenzione, Logging e Backup}
Il \textbf{Logging} è un controllo reattivo fondamentale per la \textit{Digital Forensics} e l'\textit{Incident Response}. I log devono essere raccolti, protetti da manomissioni (es. inoltrandoli a un server \textit{syslog} remoto), ruotati e analizzati (preferibilmente in modo automatizzato tramite soluzioni SIEM - \textit{Security Information and Event Management}). 

Parallelamente, rigorose policy di \textbf{Backup e Archiviazione} garantiscono la \textit{Disaster Recovery}, proteggendo l'organizzazione da guasti fisici, errori umani o attacchi distruttivi.

\section{Sicurezza in Ambiente Linux/Unix}
I sistemi \textit{Unix-like} basano storicamente la loro sicurezza su un modello di permessi granulare noto come Controllo degli Accessi Discrezionale (DAC).

\begin{itemize}
    \item \textbf{Utenti, Gruppi e ACL:} I permessi classici (suddivisi in lettura, scrittura, esecuzione per \textit{owner, group, others}) possono essere estesi tramite le \textit{Access Control Lists} (ACL) per gestire scenari di condivisione più complessi. L'autenticazione è tipicamente demandata e centralizzata tramite il framework \textbf{PAM} (\textit{Pluggable Authentication Modules}).
\end{itemize}

\section{Sicurezza in Ambiente Windows}
I sistemi operativi Microsoft Windows utilizzano un'architettura di sicurezza basata sui \textbf{Security Identifier (SID)}, stringhe univoche che identificano utenti, gruppi e computer, garantendo una gestione unificata (es. tramite \textit{Active Directory} nei domini aziendali).

\begin{itemize}
    \item \textbf{Mandatory Integrity Control (MIC):} implementa una forma derivata del modello di sicurezza di \textit{Biba}. Assegna rigidi livelli di integrità (Basso, Medio, Alto, Sistema) ai processi e agli oggetti. Questo meccanismo impedisce a codice non attendibile o isolato (come un processo del browser a integrità bassa) di modificare file di sistema eseguiti a integrità alta.
    \item \textbf{User Account Control (UAC):} Rappresenta l'implementazione del principio del minimo privilegio applicato all'interfaccia utente. Gli utenti (anche se amministratori) operano di default con un token di accesso standard e restrittivo, richiedendo un'elevazione dei privilegi esplicita e consensuale per eseguire task amministrativi o modificare impostazioni di sistema.
    \item \textbf{Registro di Sistema e Crittografia:} L'hardening in Windows passa spesso dalla configurazione massiva e centralizzata di chiavi di registro, distribuita tramite \textit{Group Policy Objects} (GPO). È inoltre fondamentale abilitare la \textit{Data at Rest Encryption} nativa (es. \textbf{BitLocker} per la cifratura del disco intero o EFS per singoli file).
\end{itemize}

\section{Sicurezza della Virtualizzazione}
La virtualizzazione astrae le risorse hardware fisiche (CPU, RAM, storage, network) fornendole in modo isolato a molteplici Macchine Virtuali (VM). L'\textbf{Hypervisor} (o \textit{Virtual Machine Monitor} - VMM) è il cuore di questa architettura, responsabile del partizionamento logico, dello \textit{scheduling} e del rigido isolamento delle VM ospiti.

\subsection{Tipologie di Hypervisor}
\begin{enumerate}
    \item \textbf{Tipo 1 (Nativo / Bare-Metal):} Installato ed eseguito direttamente sull'hardware nudo del server (es. Proxmox, Microsoft Hyper-V). Offre prestazioni nettamente superiori e una superficie di attacco minore, in quanto non dipende da un OS ospite sottostante.
    \item \textbf{Tipo 2 (Hosted):} Eseguito come una normale applicazione all'interno di un OS preesistente (es. VirtualBox, VMware Workstation). È più flessibile per l'uso \textit{client} su macchine di sviluppo, ma eredita in blocco tutte le vulnerabilità, i ritardi e l'overhead dell'OS \textit{host}.
\end{enumerate}

\begin{center}
\begin{tikzpicture}[
    box/.style={draw, rectangle, minimum height=1cm, minimum width=2.5cm, align=center},
    node distance=0.2cm,
    font=\small
]
% Type 1
\node[box, fill=gray!20, minimum width=6cm] (hw1) {Hardware Fisico};
\node[box, fill=blue!20, minimum width=6cm, above=of hw1] (hyp1) {Hypervisor (Tipo 1 / Bare-Metal)};
\node[box, fill=green!20, above=of hyp1, xshift=-1.7cm] (vm1a) {App\\OS Ospite};
\node[box, fill=green!20, above=of hyp1, xshift=1.7cm] (vm1b) {App\\OS Ospite};

% Type 2
\node[box, fill=gray!20, minimum width=6cm, right=1.5cm of hw1] (hw2) {Hardware Fisico};
\node[box, fill=yellow!20, minimum width=6cm, above=of hw2] (os2) {Sistema Operativo Host};
\node[box, fill=blue!20, minimum width=6cm, above=of os2] (hyp2) {Hypervisor (Tipo 2 / Hosted)};
\node[box, fill=green!20, above=of hyp2, xshift=-1.7cm] (vm2a) {App\\OS Ospite};
\node[box, fill=green!20, above=of hyp2, xshift=1.7cm] (vm2b) {App\\OS Ospite};

\node[above=2.8cm of hw1, font=\bfseries] {Architettura di Tipo 1};
\node[above=3.6cm of hw2, font=\bfseries] {Architettura di Tipo 2};

\end{tikzpicture}
\end{center}

\subsection{Minacce e Isolamento negli Ambienti Virtualizzati}
L'introduzione dell'hypervisor aggiunge un nuovo e critico strato architetturale:
\begin{itemize}
    \item \textbf{Traffico di Rete Virtuale:} Le VM comunicano frequentemente tra loro tramite \textit{switch} virtuali interni all'hypervisor, rendendo il traffico di rete intra-host completamente invisibile ai tradizionali IPS e firewall hardware perimetrali. È necessario implementare \textbf{Firewall Virtuali} (operanti a livello dell'hypervisor o come \textit{appliance} VM dedicate) e separare rigidamente il traffico di gestione (\textit{management}), di infrastruttura (es. la rete dedicata alla migrazione live delle VM) e il traffico applicativo degli utenti.
\end{itemize}

\begin{warnbox}[VM Escape]
Il \textbf{VM Escape} è una vulnerabilità critica e devastante in ambienti multi-tenant (come il Cloud). Si verifica quando un attaccante, partendo dal controllo completo di una VM compromessa, riesce a sfruttare un bug nell'hypervisor per eludere l'isolamento ed eseguire codice arbitrario direttamente sull'hypervisor host o all'interno delle altre VM co-ospitate sullo stesso hardware fisico.
\end{warnbox}

\section{Metasploit Framework: Penetration Testing e Assessment}
Il \textbf{Metasploit Framework} è l'infrastruttura open-source \textit{de facto} per lo sviluppo, il collaudo e l'esecuzione di codice \textit{exploit}. È utilizzato costantemente dai professionisti della sicurezza (\textit{White Hat} e \textit{Red Team}) per validare sperimentalmente l'efficacia delle misure di \textit{hardening} e condurre complessi \textit{Penetration Test}.

\begin{center}
\begin{tikzpicture}[
    box/.style={draw, rectangle, rounded corners, minimum height=1.2cm, minimum width=2.5cm, align=center},
    arrow/.style={-Stealth, thick},
    node distance=1.5cm,
    font=\small
]

\node[box, fill=blue!10] (recon) {\textbf{1. Auxiliary}\\(Ricognizione e\\Scansione)};
\node[box, fill=red!20, right=of recon] (exploit) {\textbf{2. Exploit}\\(Sfruttamento\\Vulnerabilità)};
\node[box, fill=yellow!20, right=of exploit] (payload) {\textbf{3. Payload}\\(Esecuzione\\es. Meterpreter)};
\node[box, fill=green!20, right=of payload] (post) {\textbf{4. Post}\\(Pivoting e\\Mantenimento)};

\draw[arrow] (recon) -- (exploit);
\draw[arrow] (exploit) -- (payload);
\draw[arrow] (payload) -- (post);

\end{tikzpicture}
\end{center}

\subsection{Architettura e Moduli}
Metasploit vanta un'architettura modulare ed estensibile che permette di disaccoppiare logicamente la fase di ingaggio (\textit{exploit}) dal software esecutivo (\textit{payload}). I moduli si dividono in:

\begin{warnbox}[Considerazioni Etiche e Legali (Dual-Use)]
Metasploit, al pari di molti strumenti di sicurezza, è un software a duplice uso (\textit{dual-use}). Nel contesto accademico e professionale aziendale, il suo utilizzo è strettamente normato e sorvegliato. L'esecuzione non autorizzata di moduli di exploit verso infrastrutture e server di cui non si possiede la proprietà, o per le quali non si è ottenuta un'autorizzazione scritta, formale ed esplicita (\textit{Scope of Work} / \textit{Rules of Engagement}), costituisce un \textbf{grave reato informatico} severamente punito dalle normative internazionali (es. accesso abusivo a sistema informatico).
\end{warnbox}
